\documentclass[10pt,a4paper]{amsart}
\usepackage[margin=19mm]{geometry}
\usepackage{amsmath,amssymb,mathtools}
\usepackage[scr=rsfs,cal=euler]{mathalfa}
\usepackage{xfrac}
\usepackage[unicode,hidelinks,bookmarks=false]{hyperref}
\newcommand{\ie}{i.e.\ }
\newcommand{\cf}{cf.\ }
\newcommand{\resp}{respectively\ }
\newcommand{\Subsection}{\subsection}
\providecommand{\nobreakdash}{\nobreak\hbox{-}}
\setlength{\emergencystretch}{2em}
\multlinegap0pt
\usepackage{bm}
\usepackage{bbm}
\usepackage{dsfont}
\usepackage{xspace}
\usepackage{cancel}
\usepackage{mathtools}
\usepackage{amsthm}
\makeatletter
\@ifundefined{equation}{\newtheorem{equation}{Equation}}{}
\@ifundefined{lemma}{\newtheorem{lemma}[equation]{Lemma}}{}
\@ifundefined{proposition}{\newtheorem{proposition}[equation]{Proposition}}{}
\makeatother
\newcommand{\Adm}{\operatorname{Adm}}
\newcommand{\Mod}{{\rm Mod}}
\newcommand{\N}{\mathbb{N}}
\newcommand{\Res}{{\rm Res}}
\newcommand{\cE}{\mathcal{E}}
\newcommand{\cJ}{\mathcal{J}}
\newcommand{\cM}{\mathcal{M}}
\newcommand{\divr}{{\rm div}}
\newcommand{\len}{{\rm len}}
\title[Conformal dimension and its attainment on self-similar Laaks]{Conformal dimension and its attainment on self-similar Laakso-type fractal spaces}
\author{Riku Anttila, Sylvester Eriksson-Bique, Lassi Rainio}
\date{Source: arXiv:2601.14241v1 (CC BY 4.0)}
\begin{document}
\maketitle

\noindent\emph{Source and licence.} Conformal dimension and its attainment on self-similar Laakso-type fractal spaces. Version arXiv:2601.14241v1 (CC BY 4.0), distributed under the Creative Commons Attribution 4.0 International licence, as recorded in the arXiv licence metadata for this exact version and shown on the abstract page https://arxiv.org/abs/2601.14241v1. Full rights notice: licenses/arxiv-2601.14241-CC-BY-4.0.txt.\par\smallskip

\noindent\textit{Editorial scope.} Counted unit: the Lemma whose source label is \texttt{Lemma: Optimal Density} in the article (source0.tex lines 1713--1720, source label \texttt{Lemma: Optimal Density}), delivered together with the article's own complete original proof of it (source0.tex lines 1721--1805, one \texttt{proof} environment of 85 lines). No mathematics is written, repaired or re-derived: statement and proof are the article's own text line for line. Required context retained uncounted: Lemma: Existence of Minimizers (lines 805--807); Proposition: Duality (lines 884--898); Proposition: Path lift (lines 1061--1067). Every definition, display and label of those lines is retained unchanged. Omitted from this package and not redistributed: 1--804, 808--883, 899--1060, 1068--1712, 1806--2275. Nothing inside a retained excerpt is omitted or altered: every retained body is delivered exactly as the article's lines are written, so no omission receipt of any kind accompanies this package. Citations: the retained excerpts carry no citation command at all, so no bibliography entry of the article is transcribed and no reference is redistributed. Reference adaptation: none was needed and none was made; every reference inside the delivered unit resolves inside it, so no unresolved reference or citation is produced. Printed slips kept literal: no printed word, symbol, hypothesis or number of the counted statement or of its proof was altered. Layout and class: the article's own class is \texttt{amsart} with packages including graphicx, hyperref, amsmath, amsthm, amssymb, amscd, mathtools, enumerate, verbatim, cite, esint, tikz-cd, enumitem; the wrapper is the lane's pinned \texttt{amsart} editorial wrapper at 10pt on A4, which restates the theorem-like environments the retained text needs and adds the article's own macro definitions that the retained text needs (\texttt{\textbackslash Adm}, \texttt{\textbackslash Mod}, \texttt{\textbackslash N}, \texttt{\textbackslash Res}, \texttt{\textbackslash cE}, \texttt{\textbackslash cJ}, \texttt{\textbackslash cM}, \texttt{\textbackslash divr}, \texttt{\textbackslash len}), with extra wrapper packages (bm, bbm, dsfont, xspace, cancel, mathtools). The delivered numbering is the unit's own. Assets: no retained span contains a figure environment, a graphics-inclusion command, a PDF-page inclusion, a table or a TikZ picture; no image file is copied into this package, the package declares no asset dependency and no image file is redistributed with it. Licence: version arXiv:2601.14241v1 of this article is distributed under the Creative Commons Attribution 4.0 International licence, as recorded in the arXiv licence metadata for this exact version and shown on the abstract page https://arxiv.org/abs/2601.14241v1. A full rights notice is delivered at licenses/arxiv-2601.14241-CC-BY-4.0.txt. Characters: every retained excerpt is pure ASCII, so the delivered engine, XeLaTeX with its default Latin Modern text font, reports no missing character.
\begin{lemma}\label{Lemma: Existence of Minimizers}
    Let \(\Theta\) be a non-empty family of paths on a graph \(G=(V,E,\xi)\). Then, there exists a $p$-optimal $\rho^*\in\Adm(\Theta)$. If \(p>1\), then the $p$-optimal $\rho^*$ is unique. Further,  if \(\len(\theta)\geq 2\) for all \(\theta\in\Theta\), then \(\rho(e)<1\) for all \(e\in E\).
\end{lemma}

\begin{proposition}\label{Proposition: Duality}
    Let $p\in (1,\infty)$ and $q = p/(p-1)$. Let \(G=(V,E,\xi)\) be a graph and let \(A,B\subset V\) be a pair of disjoint non-empty subsets such that $\Theta(A,B) \neq \varnothing$.
    Then
        \[
            \Mod_p(\Theta(A,B),G)^{\frac{1}{p}}\cdot\Res_p(A,B,G)^{\frac{1}{q}}=1.
        \]
    Further, if \(\cJ\) is a unit flow from \(A\) to \(B\) and \(\rho\in \Adm(\Theta(A,B))\), such that
        \[
            \mathcal{M}_p(\rho)^{\frac{1}{p}}\cdot\mathcal{E}_q(\mathcal{\cJ})^{\frac{1}{q}}=1,
        \]
    then both $\cJ$ and $\rho$ are optimal. For optimal $\cJ$ and $\rho$ we have
        \begin{equation}\label{flow_mod_vanish}
            \lvert \mathcal{F}(e) \rvert = \Mod_p(\Theta(A,B),G)^{-1}\rho(e)^{p-1}
        \end{equation}
\end{proposition}

\begin{proposition}\label{Proposition: Path lift}
    Let \(\theta=[v_0,e_1,v_1,\dots ,e_k,v_k]\) be a path in \(G_{n+m}\) and \(A,B\subset V_n\) be non-empty disjoint sets. If \(v_0\in \pi_{n+m,n}^{-1}(A)\) and \(v_k\in \pi_{n+m,n}^{-1}(B)\), then there are numbers \(0\leq k_1 < s_1 \leq k_2 < s_2\leq,\dots,\leq k_l < s_l\leq k\), vertices \(u_0,\dots ,u_l\in V_n\) and edges \(f_1,\dots ,f_l\in E_n\) so that the following conditions hold
    \begin{enumerate}
        \item \([u_0,f_1,u_1\dots,f_l,u_l]\in \Theta(A,B)\),
        \item For $i=1,\dots,l$, the sub-path \(\theta_i=[v_{k_i},e_{k_i+1},\dots ,e_{s_i}, v_{s_i}]\) of \(\theta\) is a path from  \(v_{k_i}\in\pi_{n+m,n}^{-1}(u_i)\) to \(v_{s_i}\in\pi_{n+m,n}^{-1}(u_{i+1})\) and is contained in \(f_i\cdot G_m\).
    \end{enumerate}
\end{proposition}

\begin{lemma}\label{Lemma: Optimal Density}
    For every \(n\in\N\) and \(p\in(1,\infty)\) the unique \(p\)-minimal \(\Theta^{(n)}\)-admissible density \(\rho_n\colon E_n\to[0,1]\) is symmetric, satisfies \(\rho_n(e)<1\) for all \(e\in E_n\), and admits a multiplicative cascade representation
    
        \begin{equation}\label{Equation: Optimal density is a cascade}
            \rho_n(e)=\prod_{i=1}^n\rho_1(T^{-1}(e)_i)\quad\text{for all }e\in E_n,
        \end{equation}
    where \(\rho_1\colon E_1\to [0,1)\) is the \(p\)-minimal \(\Theta^{(1)}\)-admissible density. Furthermore, the mapping \(t\mapsto\Mod_t(\Theta^{(n)},G_n)\) is continuous and strictly decreasing for all \(t\in[1,\infty)\).
\end{lemma}

\begin{proof}
    
    Fix \(p\in(1,\infty)\). Since the IGS is non-degenerate, we have that \(\len(\theta)\geq 2\) for all \(\theta\in\Theta^{(n)}\). It then follows directly from Lemma \ref{Lemma: Existence of Minimizers} that the unique \(p\)-minimal \(\Theta^{(n)}\)-admissible density \(\rho_n\colon E_n\to[0,1]\) satisfies \(\rho_n(e)<1\) for all \(e\in E_n\), and that the mapping \(t\mapsto\Mod_t(\Theta^{(n)},G_n)\) is continuous and strictly decreasing for all \(t\in[1,\infty)\). Since the IGS is symmetric, we have
        \[
            \cM_p(\rho_n\circ \eta)=\cM_p(\rho_n)=\Mod_p(\Theta^{(n)},G_n).
        \]
    Since the minimizer $\rho_n$ is unique, we necessarily have \(\rho_n\circ \eta=\rho_n\). 
    
    We now verify \eqref{Equation: Optimal density is a cascade} for each \(n\in\N\). First, by applying Proposition \ref{Proposition: Path lift} inductively, we see that
        \[
            \min_{\theta\in\Theta^{(n)}}L_{\rho_n}(\theta)\geq\big(\min_{\theta\in\Theta^{(1)}}L_{\rho}(\theta)\big)^n\geq1.
        \]
    Therefore, \(\rho_n\) is \(\Theta^{(n)}\)-admissible. Due to Proposition \ref{Proposition: Duality} and the uniqueness of \(\rho_n\), it is enough to show that there exists a unit flow \(\cJ_n\) from \(I^{(n)}_-\) to \(I^{(n)}_+\) such that, if \(\rho_n\) satisfies equation \eqref{Equation: Optimal density is a cascade}, then 
        \[
            \cM_p(\rho_n)^\frac{1}{p}\cdot \cE_q(\cJ_n)^\frac{1}{q}=1.
        \]
    Since \(\cM_p(\rho_1)=\Mod_p(\Theta^{(1)},G_1)\), there exists by Proposition \ref{Proposition: Duality} a unit flow \(\cJ_1\) from \(I_-\) to \(I_+\) so that \(\cM_p(\rho_1)^\frac{1}{p}\cdot\cE_q(\cJ_1)^\frac{1}{q}=1\). Note, that since \(\rho_1\circ \eta = \rho_1\), it follows that \(\cJ_1(\eta(z),\eta(e))=-\cJ_1(z,e)\) for every \(z\in e\in E_1\). For every \(n\in\N\), we define \(\cJ_{n+1}\) as follows: For every \(z\in e\in E_1\) and \(f\in E_n\) we set
        \[
            \cJ_{n+1}([z,f],(e,f)):=\cJ_{n}(f^-,f)\cJ_1(z,e).
        \]
    We will show by induction, that \(\cJ_{n+1}\) is a unit flow for from \(I_-^{(n+1)}\) to \(I_+^{(n+1)}\) for every \(n\in \N\). To that end, assume that \(\cJ_n\) is a unit flow. Since \(\cJ_1\) is unit flow, we have for all \(e\in E_1\) and \(f\in E_n\) that
        \[
        \begin{split}
            \cJ_{n+1}([e^+,f],(e,f))=\cJ_{n}(f^-,f)\cJ_1(e^+,e)
            &=-\cJ_{n}(f^-,f)\cJ_1(e^-,e)\\&=-\cJ_{n+1}([e^+,f],(e,f)).
        \end{split}
        \]
    Now, let \(v=[z,\hat{e}]\in e\in E_{n+1}\) for some \(z\in V_1\) and \(\hat{e}\in E_n\), and suppose that \(z\not\in I_-\cup I_+\). The edges with \(v\) as an endpoint are all contained in \(\hat{e}\cdot G_1\), and since \(\cJ_1\) is a flow, we have
        \[
        \begin{split}
            \divr(\cJ_{n+1})(v)=\sum_{\{f\mid v\in f\}}\cJ_{n+1}(v,f)&=\sum_{\{f\mid z\in f\}}\cJ_n(\hat{e}^-,\hat{e})\cJ_1(z,f)\\
            &=\cJ_n(\hat{e}^-,\hat{e})\cdot\divr(\cJ_1)(z) =0.
        \end{split}
        \]
    Suppose next that \(z\in I_-\cup I_+\). Then, \(v\in\pi_{n+1}^{-1}(\hat{v})\) for some \(\hat{v}=\pi_{n+1}(v)\in V_n\). Since the IGS is doubling, for every edge \(\hat{f}\in E_n\) with \(\hat{v}\) as an endpoint the set \(\hat{f}\cdot E_1\subset E_{n+1}\) contains exactly one edge \(f\) with \(v\) as an endpoint, and by construction no other edges with \(v\) as an endpoint can exist. Since the IGS is symmetric, there exists a vertex \(x\in I_-\) so that \(z\in\{x,\eta(x)\}\) and hence, for each of the edges \(\hat{f}\) with \(\hat{v}\) as an endpoint, the vertex \(v\) satisfies:
    \[
        v=[z,\hat{e}]=
        \begin{cases}
            [x,\hat{f}]&\text{if }\hat{v}=\hat{f}^-,\\
            [\eta(x),\hat{f}]&\text{if }\hat{v}=\hat{f}^+.
        \end{cases}
    \]
    Therefore, we have that
    \[
        \cJ_{n+1}(v,f):=
        \begin{cases}
            \cJ_n(\hat{v},\hat{f})\cJ_1(x,\mathfrak{e}(x))  &\text{if }\hat{v}=\hat{f}^-,\\
            -\cJ_n(\hat{v},\hat{f})\cJ_1(\eta(x),\mathfrak{e}(\eta(x))) &\text{if }\hat{v}=\hat{f}^+,
        \end{cases}
    \]
    and since \(\mathfrak{e}(\eta(x))=\eta(\mathfrak{e}(x))\) and \(\cJ_1(x,\mathfrak{e}(x))=-\cJ_1(\eta(x),\eta(\mathfrak{e}(x)))\), it then follows that
    \[
        \cJ_{n+1}(v,f)=\cJ_n(\hat{v},\hat{f})\cJ_1(x,\mathfrak{e}(x)),
    \]
    from which we may compute
    \begin{equation}\label{Equation: Divergence Lift}
    \begin{split}
        \divr(\cJ_{n+1})(v)=\sum_{\{f\mid v\in f\}}\cJ_{n+1}(v,f)&=\sum_{\{\hat{f}\mid \hat{v}\in\hat{f}\}}\cJ_n(\hat{v},\hat{f}) \cJ_1(x,\mathfrak{e}(x))\\
        &=\cJ_1(x,\mathfrak{e}(x))\cdot\divr(\cJ_n)(\hat{v}).
    \end{split}
    \end{equation}
    If \(v\not\in I_-^{(n+1)}\cup I_+^{(n+1)}\), then \(\hat{v}\not\in I_-^{(n)}\cup I_+^{(n)}\), and consequently \(\divr(\cJ_n)(\hat{v})=0\), and so we have by \eqref{Equation: Divergence Lift} that \(\divr(\cJ_{n+1})(v)=0\). Hence, \(\cJ_{n+1}\) is a flow. To see that it is a unit flow, we compute inductively the total flow:
    \[
    \begin{split}
        \sum_{v\in I_-^{(n+1)}}\divr(\cJ_{n+1})(v)&\stackrel{\eqref{Equation: Divergence Lift}}{=}\sum_{\hat{v}\in I_-^{(n)}}\sum_{x\in I_-^{(1)}}\cJ_1(x,\mathfrak{e}(x))\cdot \divr(\cJ_n)(\hat{v})\\
        &\stackrel{\phantom{\eqref{Equation: Divergence Lift}}}{=}\sum_{\hat{v}\in I_-^{(n)}}\divr(\cJ_n)(\hat{v})\sum_{x\in I_-^{(1)}}\cJ_1(x,\mathfrak{e}(x))\\
        &\stackrel{\phantom{\eqref{Equation: Divergence Lift}}}{=}\sum_{\hat{v}\in I_-^{(n)}}\divr(\cJ_n)(\hat{v})=1
    \end{split}
    \]
    Having constructed the unit flow \(\cJ_n\) from \(I^{(n)}_-\) to \(I^{(n)}_+\) for every \(n\in\N\), we only need show that if \(\rho_n\) satisfies equation \eqref{Equation: Optimal density is a cascade} then \(\cM_p(\rho_n)^{1/p} \cdot \cE_q(\cJ_n)^{1/q}=1\). To that end, first note that for every \(e\in E_{n}\) we have by construction that
    \begin{equation}\label{Equation: Optimal flow is a cascade}
        \lvert \cJ_{n}(e)\rvert=\prod_{i=1}^{n}\lvert \cJ_1(T^{-1}(e)_i)\rvert.
    \end{equation}
    Then, if \(\rho_n\) satisfies equation \eqref{Equation: Optimal density is a cascade}, we may compute:
    \begin{equation}\label{Equation: Energy and mass cascade}
    \begin{split}
        \cE_q(\cJ_n)&\stackrel{\eqref{Equation: Optimal flow is a cascade}}{=}\sum_{e\in E_n}\prod_{i=1}^{n}\lvert \cJ_1(T^{-1}(e)_i)\rvert^q=\prod_{i=1}^{n}\sum_{e\in E_1}\lvert\cJ_1(e)\rvert^q=\cE_q(\cJ_1)^n,\\
        \cM_p(\rho_n)&\stackrel{\eqref{Equation: Optimal density is a cascade}}{=}\sum_{e\in E_n}\prod_{i=1}^n\lvert \rho_1(T^{-1}(e)_i)\rvert^p=\prod_{i=1}^n\sum_{e\in E_1}\lvert\rho_n(e)\rvert^p=\cM_p(\rho_1)^n.
    \end{split}
    \end{equation}
    Since \(\cM_p(\rho_1)^{1/p} \cdot \cE_q(\cJ_1)^{1/q}=1\), it then follows that
    \[
        \cM_p(\rho_n)^\frac{1}{p}\cdot \cE_q(\cJ_n)^\frac{1}{q}\stackrel{\eqref{Equation: Energy and mass cascade}}{=}\cM_p(\rho_1)^\frac{n}{p}\cdot \cE_q(\cJ_1)^\frac{n}{q}=1.
    \]  
\end{proof}

\end{document}
